\section{Additivity of native model potentials and exact FS functions}
\label{sec:native-tc-nef-fs-tensor-potential}

Let $X$ be integral, $\pi\colon\mathcal X\to\mathbf A^1$ a test configuration of $X$, and
$(M,\alpha,e)$, $(H,\beta,f)$ linearized lines with generic lines $L,K$.  Write
$\varphi_M$, $\varphi_H$ for their native model potentials at the fixed Gauss radius, and
use the linearizations of tensor products and powers of
Lemma~\ref{lem:native-tc-nef-fs-whole-tensor-comparisons}.

\begin{theorem}[Native norm factors multiply {\cite[\S1.2]{BJ22}}]
\label{thm:native-tc-nef-fs-tensor-native-factor}
\leanok
\lean{EnvelopeProperty.NativeTestConfiguration.TestConfiguration.gaussNativeNormFactor_tensor_comparisons}
\uses{lem:native-tc-nef-fs-whole-tensor-comparisons,
lem:native-tc-nef-fs-product-tensor-coherence,
lem:native-tc-nef-fs-induced-residue-tensor-frame,
lem:native-tc-nef-fs-reference-support-tensor-frame,
lem:native-tc-nef-fs-native-factor-arbitrary-center,
thm:native-tc-intrinsic-factor-original-frame-independence,
thm:actual-original-common-affine-line-frames,
lem:actual-center-generic-tensor-frame,
thm:fc-original-ample-tensor-ample,
thm:fc-actual-power-frame-transition-units}
If $L$ and $K$ are ample, then at every point and radius the native norm factor of
$M\otimes H$ is the product of the native norm factors of $M$ and of $H$.
\end{theorem}

\begin{proof}
Since $L$ and $K$ are invertible, there is an affine open $U$ containing the centre of the
given point on which both $L$ and $K$ are trivial; fix frames $e_L,e_K$ there and use the
tensor frame $e_L\otimes e_K$ for $L\otimes K$.  At the valuation centre of the Gauss point
in the total space, use the tensor of the two chosen centre frames of $M$ and $H$; its
generic frame is the tensor of their generic frames
(Lemma~\ref{lem:actual-center-generic-tensor-frame}).

The line $L\otimes K$ is ample because $X$ is integral
(Theorem~\ref{thm:fc-original-ample-tensor-ample}).  Hence, by
Lemma~\ref{lem:native-tc-nef-fs-native-factor-arbitrary-center}, the native factor of
$M\otimes H$ is the valuation of the transition unit between the reference frame induced
by $e_L\otimes e_K$ and the model frame induced by the tensor centre frame.  By
Theorem~\ref{thm:native-tc-intrinsic-factor-original-frame-independence} the native factors
of $M$ and $H$ are computed in the same way with $e_L$, resp.\ $e_K$, on the same open $U$.

The product comparison of the tensor data is the tensor of the two product comparisons,
conjugated by the pullback tensorators
(Lemma~\ref{lem:native-tc-nef-fs-product-tensor-coherence}); induced residue frames and
reference support frames commute with tensor products
(Lemmas~\ref{lem:native-tc-nef-fs-induced-residue-tensor-frame}
and~\ref{lem:native-tc-nef-fs-reference-support-tensor-frame}).  Therefore both the
reference frame and the model frame of $M\otimes H$ are the same inverse tensorator
followed by the tensor frame of the frames of $M$ and $H$.  A common preceding isomorphism
cancels in a transition unit, and the transition unit of two tensor frames is the product
of the transition units.  The residue valuation is multiplicative.
\end{proof}

\begin{theorem}[Native model potentials add]
\label{thm:native-tc-nef-fs-tensor-native-potential}
\leanok
\lean{EnvelopeProperty.NativeTestConfiguration.TestConfiguration.gaussNativeModelPotential_tensor_comparisons}
\uses{thm:native-tc-nef-fs-tensor-native-factor, def:native-tc-original-model-potential}
If $L$ and $K$ are ample, the native model potential of $M\otimes H$ is
$\varphi_M+\varphi_H$.
\end{theorem}

\begin{proof}
The native potential is minus the logarithm of the positive native factor.  Apply
Theorem~\ref{thm:native-tc-nef-fs-tensor-native-factor} and $\log(cd)=\log c+\log d$ for
$c,d>0$.
\end{proof}

\begin{theorem}[Native potential of a power-tensor line]
\label{thm:native-tc-nef-fs-power-tensor-native-potential}
\leanok
\lean{EnvelopeProperty.NativeTestConfiguration.TestConfiguration.gaussNativeModelPotential_powerTensor_comparisons}
\uses{thm:native-tc-nef-fs-tensor-native-potential, thm:native-tc-nef-fs-power-potential}
If $L^{\otimes b}$ and $K^{\otimes a}$ are ample, the native model potential of
$M^{\otimes b}\otimes H^{\otimes a}$ is $b\varphi_M+a\varphi_H$.
\end{theorem}

\begin{proof}
Apply Theorem~\ref{thm:native-tc-nef-fs-tensor-native-potential} to the linearized lines
$M^{\otimes b}$ and $H^{\otimes a}$, whose generic lines $L^{\otimes b}$, $K^{\otimes a}$
are ample, and then Theorem~\ref{thm:native-tc-nef-fs-power-potential} to each factor.
\end{proof}

\begin{theorem}[Exact FS functions from a centrally nef line {\cite[Lemma~5.5]{FHX}}]
\label{thm:native-tc-nef-fs-actual-potential-sum-finite-fs}
\leanok
\lean{EnvelopeProperty.NativeTestConfiguration.TestConfiguration.isFinFS_power_tensor_potential_sum_of_central_nef}
\uses{thm:native-tc-nef-fs-perturbed-model-exact,
thm:native-tc-nef-fs-power-tensor-native-potential}
Assume moreover that $k$ is infinite, $\mathcal X$ is integral and dominant over
$\mathbf A^1$, $M|_{\mathcal X_0}$ is nef and $H|_{\mathcal X_0}$ is ample.  If $a>0$ and
$L^{\otimes b}$, $K^{\otimes a}$ are ample, then $b\varphi_M+a\varphi_H$ is a finite
rational Fubini--Study function of class $1$ for $L^{\otimes b}\otimes K^{\otimes a}$.
\end{theorem}

\begin{proof}
By Theorem~\ref{thm:native-tc-nef-fs-perturbed-model-exact} the native model potential of
$M^{\otimes b}\otimes H^{\otimes a}$ is such a function, and by
Theorem~\ref{thm:native-tc-nef-fs-power-tensor-native-potential} it equals
$b\varphi_M+a\varphi_H$ at every point.
\end{proof}

\begin{corollary}[Division by the nef exponent {\cite[(2.2)]{FHX}}]
\label{cor:native-tc-nef-fs-potential-add-div}
\leanok
\lean{EnvelopeProperty.NativeTestConfiguration.TestConfiguration.isFinFS_potential_add_div_of_central_nef}
\uses{thm:native-tc-nef-fs-actual-potential-sum-finite-fs,
lem:native-tc-nef-fs-ray-denominator-division}
Under the hypotheses of Theorem~\ref{thm:native-tc-nef-fs-actual-potential-sum-finite-fs}
and $b>0$, the function $\varphi_M+\frac ab\varphi_H$ is a finite rational Fubini--Study
function of class $1/b$ for $L^{\otimes b}\otimes K^{\otimes a}$, defined by the same
sections and weights.
\end{corollary}

\begin{proof}
Divide the realization of Theorem~\ref{thm:native-tc-nef-fs-actual-potential-sum-finite-fs}
by $b$ (Lemma~\ref{lem:native-tc-nef-fs-ray-denominator-division}); the resulting
function is $b^{-1}(b\varphi_M+a\varphi_H)=\varphi_M+\frac ab\varphi_H$.
\end{proof}

\begin{remark}
In the application to the local positivity lemma of Fang--Han--Xia, $\varphi_M$ is
(up to the normalizations $r\varphi=\varphi_M+c$) the given locally convex function and
$H$ is an auxiliary line ample on the central fibre.  Letting $b\to\infty$ for fixed $a$
then exhibits $\varphi_M$ as a uniform limit of Fubini--Study functions, provided
$\varphi_H$ is bounded.  The construction of $H$, the bound on $\varphi_H$, and the
passage from these native functions to Fubini--Study functions on $X^{\mathrm{an}}$ are
not part of this chapter.
\end{remark}
